% !TeX root = ../main.tex

%------------------------------- 1
\chapter{引言}
%引言的写法：1.~问题的提出；2.~选题背景及意义；3.~文献综述；4.~研究方法；5.~论文结构安排。

  人类通过科学探索和研究，不断加深对物质世界的理解，并推动了科学技术的发展。这一过程在古代至今持续进行，为我们揭示了自然界的奥秘。
  在古代，人类对物质世界的理解主要基于直观和经验。例如，古希腊的德谟克利特的原子论思想，以及古代中国的阴阳五行学说。这种理论简单的为物质的复杂性和多样性提供了初步的解释。
  1687年7月，艾萨克·牛顿出版了《自然哲学的数学原理》一书，标志着经典力学体系的正式建立。
  1897年，卡文迪许实验室的约瑟夫·汤姆逊发现了电子。
  1919年，卢瑟福使用α粒子（氦核）轰击氮原子，发现质子。
  1932年，卡文迪许实验室的詹姆斯·查德威克发现中子。
  1935年，汤川秀树提出重要的理论，解释了原子核内质子和中子之间的强相互作用。
  1900年，普朗克公式被提出可以通过引入振子能量量子化的概念导出。量子力学由此诞生。
  1964年，盖尔曼和乔治·茨威格分别独立提出夸克模型。
  1965年，南部阳一郎等人分别独立提出夸克应当具有一个额外的SU(3)规范对称的自由度，这种自由度后来被称作色荷。
  1969年，费曼提出了部分子模型（parton model）。
  1979年9月，由丁肇中领导的高能物理实验小组，通过实验找到了胶子存在的证据。
  2012年，科学家们宣布在大型强子对撞机（LHC）上发现了希格斯玻色子（上帝粒子）。
  从波尔的原子模型，到夸克模型，再到标准模型，这些理论为我们揭示了物质的微观世界和更深层次的秘密。
  
	组成物质的基本粒子包括夸克（Quarks）、轻子（Leptons）、玻色子（Bosons）、引力子（Graviton）。这些基本粒子之间存在多种相互作用，比如由量子色动力学（Quantum Chromodynamics，简称QCD）描述的强相互作用，它作用于带有色荷的夸克和胶子（gluon），通过胶子来传递；
	由量子电动力学（Quantum electrodynamics，简称QED）描述的电磁相互作用，它作用于带有电荷的粒子，例如电子和光子，通过光子来传递；
	弱相互作用作用于带有弱荷的粒子，例如带电轻子和中微子，通过W玻色子和Z玻色来传递相互作用；
	由广义相对论描述的引力相互作用，它作用于物质和能量，包括质量引力和时空的弯曲，由尚未被实验观测到的引力子传递。
	高能核物理学是物理学的一个分支，探讨原子核及其成分，如质子和中子在极高能量下的行为。现代物理学的基础理论之一即量子色动力学。它起源于20世纪70年代早期，在理解基本粒子物理和强子的性质方面起着重要作用。
	高能重离子（heavy ion colision）碰撞实验提供了研究量子色动力学的独特机会。相对论重离子碰撞物理一直是对基本粒子和宇宙起源做深入研究的重要前沿领域，是研究高能重离子（如铅离子或金离子）在接近光速下发生碰撞时所涉及的物理现象和过程的学科领域。这种高能重离子的碰撞能够创造出极高的温度和能量密度，使得物质处于一种称为夸克胶子等离子体（Quark-Gluon Plasma，简称QGP）的状态。QGP是一种类似于宇宙大爆炸初期的物质形态，其中夸克和胶子分离，并自由存在于热密度极高的介质中。这种高能密度下的物质态是研究QCD的重要实验环境之一。
	对该领域的研究为我们了解宇宙的起源和性质提供了关键信息，深化了我们对物质世界的理解。

%------------------------------- 1.1

\section{标准模型和量子色动力学}
\label{sec:SM}
	实验学家已经确定了已知物质粒子的行为。有了这些知识，我们可以确定粒子的相互作用如何影响宇宙中的重要过程。
	物理学表明，基本粒子扮演着许多角色。夸克和轻子是物质的基本组成部分，它们结合成质子、中子和原子，构成了地球上生命的基础。所谓的规范粒子产生了各种力，包括电、磁和重力。
	
	标准模型（Standard Model，简称SM）~\parencite{D0:1995jca,Weinberg:1973un,quigg1985gauge}是描述基本粒子及其相互作用的理论框架。
	自然界中的基本相互作用有四种即引相互作用、电磁相互作用、弱相互作用和强相互作用。标准模型能够描述除引力以外的三种基本相互作用以及所有基本粒子，包括夸克、轻子、规范玻色子和希格斯玻色子~\parencite{weinberg2004making}。
	不同的相互作用通过费米子组分之间的特征玻色子（具有整数自旋的粒子）的交换来用量子语言描述。其中强相互作用负责将夸克束缚在中子和质子中，以及将中子和质子束缚在原子核中。夸克之间的相互作用由无质量的粒子胶子传递。
	胶子是质量为零且自旋为$1$的玻色子，类似于光子。然而，由于胶子的色荷相互作用，胶子之间可以有多达$N_c^2 - 1 = 8$种颜色。
	\iffalse
	电磁相互作用负责几乎所有的超核物理现象，例如电子与核之间的束缚态，即原子和分子。这些相互作用通过光子的交换来传递。
	弱相互作用的典型例子是核$\beta$衰变的缓慢过程，涉及放射性核发射出一个电子和中微子。弱相互作用的传递者是$W^{\pm}$和$Z^0$玻色子，它们的质量约为质子质量的100倍。
	引相互作用作用于所有类型的粒子。在粒子物理实验的尺度上，引力是所有基本相互作用中最弱的。
	\fi
	自20世纪中叶以来，高能物理的实验成果都与标准模型的预测相吻合。
	在对高能核碰撞物理领域的研究中，标准模型发挥着至关重要的作用。
	
%	标准模型中基本粒子包括：
	如图~ref{fig:StandardModel}所示，根据这个模型，所有物质都是由一小组基本自旋1/2的费米子所构成：六种夸克和六种轻子。
	夸克是自旋$1/2$的费米子，携带分数电荷$+2/3e$或$-1/3e$。夸克的类型即“味道”通过符号来表示：$u$代表“上夸克"，$d$代表下夸克、$s$代表奇异夸克、$c$代表粲夸克、$t$代表顶夸克、$b$代表底夸克。夸克还有一个额外的自由度：每种夸克味道有三种不同的颜色。
	轻子携带整数的电荷。众所周知电子$e$带有单位负电荷，而其它带电的轻子包括$\mu$介子和$\tau$介子。电中性轻子被称为中微子，用符号$\nu$表示。
	每种不同“味道”的中微子都与各自“味道”的带电轻子相对应，比如$\nu_e$、$\nu_\mu$、$\nu_\tau$。
	
	在正常的核物质中，夸克被束缚在强子（如质子和中子）中，不能单独存在。夸克之间的强相互作用具有一种特殊性质，即它们只能以如$uud$的组合形式存在。
	质子和中子由最轻的u夸克和d夸克组成，每次都是三个：质子由uud组成，中子由ddu组成。而当前宇宙的常见物质是稳定的粒子，即电子$e$以及$u$和$d$夸克。
	更重的夸克$s$、$c$、$b$和$t$也会组合成类似于质子和中子的粒子，但它们更加重，并且不稳定且迅速衰变（通常在$10^{-13}$秒内衰变为$u$、$d$的组合）。只有在人造加速器的非常高能碰撞中，或者在宇宙射线中自然生成时，才能观察到这些重而不稳定的粒子。
	
  \hspace*{\fill}
	\begin{figure}[htbp]
	\centering
	\includegraphics[width=3.8in]{StandardModel}
	\caption{标准模型的基本粒子，\uppercase\expandafter{\romannumeral1}、\uppercase\expandafter{\romannumeral2}、\uppercase\expandafter{\romannumeral3}分别代表三代夸克轻子，第四列表示规范波色子，第五列是希格斯粒子。}
	\label{fig:StandardModel}
	\end{figure}

% 标准模型概述：：基本粒子分类和特性；强相互作用、弱相互作用和电磁相互作用的描述；标准模型的拉格朗日量和规范对称性
%
% 介绍标准模型的基本框架和理论基础。
% 讨论基本粒子的分类和特性，包括夸克、轻子、质量载体粒子和希格斯玻色子。
% 阐述强相互作用、弱相互作用和电磁相互作用的描述，并介绍相应的规范对称性。


% 量子色动力学（QCD）：：QCD的起源和发展；夸克和胶子的描述；色荷和色荷传递的胶子
%
% 详细介绍量子色动力学的起源、发展和基本原理。
% 解释夸克和胶子的描述，包括它们的色荷和色荷传递的胶子。
% 阐述QCD的拉格朗日量和基本方程，包括夸克和胶子场的运动方程。


% QCD的拉格朗日量和基本方程：：QCD拉格朗日量的形式；夸克和胶子场的运动方程；渐进自由性和非渐进性质


% QCD在高能核物理中的应用：：强子对撞机和实验观测；QCD在高能核物理中的重要性；QCD计算方法和技术
%
% 探讨QCD在高能核物理中的重要性和应用。
% 介绍强子对撞机实验和相关实验观测。
% 讨论QCD在高能核物理中的计算方法和技术，包括对截面、衰变和散射等过程的计算。

% 小结：：对本章内容进行总结和概括，强调标准模型和量子色动力学在高能核物理研究中的关键性；引导读者进入下一章的主题，为后续章节的讨论奠定基础。

	%以量子场论为基础，采用了规范场论的数学框架，引入了SU(3)规范对称性。在量子色动力学中，夸克和胶子被赋予色荷（类似于电荷）并相互作用，通过交换胶子相互作用。

	
	% 总的来说，量子色动力学的背景是在探索强相互作用时，物理学家们发现强相互作用具有复杂的非线性性质，需要一种新的理论来描述和解释。量子色动力学的提出以及随后的发展为我们揭示了强相互作用的本质，为研究强子物质、基本粒子物理和宇宙学等领域的问题提供了重要的工具和框架。

	%QCD的数学形式非常复杂，涉及到非线性、非阿贝尔规范场理论的计算，因此对于高能量和高密度的强子物质，如夸克胶子等离子体，QCD的研究面临一些挑战。

  %通过QCD，我们可以研究强相互作用的许多方面，如粒子的结构、强子的性质、夸克-胶子相变等。QCD的理论预测可以通过计算机模拟、实验观测和现象学模型进行验证和研究。

\iffalse

% 坚果云/Paper10_Book/Basice pQCD/ Prologue

随着20世纪70年代非阿贝尔规范理论中的渐进自由的发现，人们逐渐认识到部分子实际上就是在夸克和胶子通过胶子耦合到色荷的量子理论中出现的粒子。渐进自由性允许我们将夸克在短距离和短时间内视为自由粒子，从而与部分子作为强子非相互作用组成部分的概念相联系。很快，人们意识到可以将许多涉及大动量转移的高能过程分解为两个部分：一个部分仅涉及硬相互作用，可以使用微扰QCD进行计算；另一个部分则需要详细的非微扰信息来描述强子由夸克和胶子组成的方式。

% 坚果云/Paper10_Book/Into_RHIC/ 1.1.1 The Quark Gluon Plasma

预期情况是，实验室中通过重离子反应可以在相对较低的能量水平（约为10-100 A·GeV实验室入射能量）产生足以形成夸克胶子等离子体（QGP）所需的高能密度。当创造出这些条件时，反应区域将存在大量夸克和胶子。这是实验室中产生自由夸克和胶子的唯一可能性，而且仅在一个小体积（约为10-100 fm$^3$）内的短时间（2-10 fm/c）内实现。
%                         .../10.1 QGP_Introduction
量子色动力学（QCD）是描述高能下强子的基本场论。然而，为了建立量子色动力学与重离子实验之间的联系，需要进行广泛的理论和实验工作。

在极端相对论的能量下，两个核的碰撞将释放能量到一个比基本强子尺度1 fm大得多且寿命更长的时空区域。这样的碰撞为我们提供了研究热而密集的强子物质的机会，其条件类似于早期宇宙的情况。这项研究的主要目标包括确定高密度下强相互作用物质的性质、夸克-强子相变的本质、验证自然界中夸克胶子等离子体的存在，并确定由格点量子色动力学计算和现象学模型预测的夸克-强子相变参数。

研究极端密度下物质的性质，为核物理和粒子物理之间建立了联系。它是对核物质性质研究的自然延伸，受到粒子物理中揭示的强相互作用动力学的启发。这项研究提供了一种实验工具，可用于研究高能量下的多体系统，超越了粒子物理研究的可能性。
%                         .../10.1.1 Theoretical expectations
在极端密度下，可能会发生解禁相变，因为许多色荷的存在会屏蔽给定的夸克反夸克（$q\bar{q}$）或夸克三重态（$qqq$）系统成员之间的束缚势能。在早期宇宙中，这个相变可能在大爆炸后的几十微秒内发生。重离子实验为实验室研究夸克-强子相变和原始夸克胶子等离子体的性质提供了工具。在理论方面，量子色动力学（QCD）描述了包括相变时的临界行为在内的强相互作用热力学。
% QCD Predictions
在量子色动力学中，已经对过渡的参数进行了研究，包括能量密度、温度、重子密度和屏蔽长度。这项研究采用了格点形式的计算机模拟和各种现象学方法。

在格点量子色动力学（lattice QCD）的计算机模拟中，人们尝试从第一原理出发计算相关物理量，而无需依赖任何简化的物理假设。这些结果为我们提供了对于零重子数密度条件下强相互作用物质临界行为的相当良好整体理解($n_B = 0$)。

  \begin{itemize}
	\item 12种基本费米子以及对应的反粒子，根据粒子间相互作用的不同，可以分为六味夸克（顶夸克、上夸克、下夸克、奇异夸克、祭夸克和底夸克）以及六种轻子（电子、$\mu$子、$\tau$子以及对应的中微子）。
	\item 规范玻色子：传递电磁相互作用的光子、传递弱相互作用的$W^+$,$W^-$和$Z^0$玻色子以及传递强相互作用的8种带”色” 胶子。粒子间的相互作用通过规范玻色子的交换来描述。
	\item 希格斯粒子：希格斯机制被认为是产生标准模型中所有基本粒子质量的原因。
\end{itemize}

标准模型中规范玻色子的规范变换可以用称为“规范群”的酉群去描述 (unitarygrouD），标准模型可以作为一个基于规范对称群 $SU(3)_c \otimes SU(2)_L \otimes SU(1)_Y$ 的量子场论。其中 $SU(2)_L \otimes SU(1)_Y$ 是电弱相互作用对称群，$SU(3)_c$ 是强相互作用对称群。
量子色动力学-$SU(3)$ 色规范群下非阿贝尔规范场论，被公认为夸克层次上最为成功的强相互作用理论。

\fi


%------------------------------------------------------------------------------------------------------------------------------------------------
%------------------------------------------------------------------------------------------------------------------------------------------------
	在上世纪70年代初，关于质子的深度非弹性电子碰撞实验揭示了核子的内部结构，核子由夸克和胶子构成。用于描述这些粒子的场论被称为量子色动力学。该理论得出一个结论：单个自由夸克或胶子无法在我们的实验室中（所谓的物理真空）进行研究或观测，因为它们受到强相互作用的束缚。然而，在我们的物理世界中，所有可观测的粒子都是无色或色中性的。另一方面，高能密度和温度下，夸克和胶子之间的强相互作用减弱，有色物质能够传播更远的距离。这也可能是早期宇宙的情况。

	量子色动力学是标准模型的一个重要组成部分，是基于规范场论描述强相互作用的理论。它包含核物理的各个方面，从夸克和胶子的行为到核子和核的性质~\parencite{Achenbach:2023pba}。
	量子色动力学是非阿贝尔规范场论的一个成功运用，它所对应是非阿贝尔规范群的 $SU(3)$~\parencite{Gross:1973id,Politzer:1973fx}，群量子数被称为“色荷”。每一种夸克的三种颜色对应着 $SU(3)$ 群的基本表示。胶子作为强相互作用力的传递者，有八种，对应着 $SU(3)$ 群的伴随表示。这个理论的动力学完全由它的 $SU(3)$ 规范对称群决定。
	
	夸克和胶子是强子的组成部分。强子内的部分子可以由三类粒子组成：一类称为价夸克，它们的数目和味是确定的，并随不同强子而不同，价夸克决定强子的性质；第二类称为海夸克，它们的数目和味是不确定的，但其总和的味性质和真空相同；第三类称为胶子，它们的数目不定，其味性质和真空相同，起传递色相互作用的作用。
	普通强子（如质子、中子、$\pi$、$\Delta$等）仅包含$u$和$d$夸克，奇异强子（如$\Lambda$、$\Omega$等）还包含$s$夸克，依此类推。

	对于质量为$m$的夸克，规范不変的拉格朗日密度为：
  \begin{align}
    \mathcal{L}_{\mathrm{QCD}}=\bar{q}^a \left(i \gamma^{\mu} D_{\mu} - m \delta_{\alpha \beta}\right) q^{\beta} - \frac{1}{4} F_{\mu \nu}^a F_a^{\mu \nu}
  \label{eq:lagran}
  \end{align}
 	其中$q^a$为代表三重态夸克场，$D_{\mu} \equiv \partial_{mu} + igt^a A_{\mu^a}$ 为作用在$q^a$的协变微商，$g$为QCD无量纲的耦合常数，$\alpha, \beta$ 代表有三种色荷的夸克，$t^a$ 是 $SU(3)$ 群的八个生成元，满足	$ {\left[t^a, t^b\right]=i f^{a b c} t^c  F_{\mu \nu}^a}$。
	$F_{\mu \nu}^a$ 是胶子场的场强张量：
	\begin{equation}
			F_{\mu \nu}^a=\partial_\mu A_\nu^a-\partial_\nu A_\mu^a-g f_{a b c} A_\mu^b A_\nu^c
	\end{equation}
	其中$ A_\mu^a$代表八重态的胶子场（a取值为1到8），$ f_{a b c}$ 为 $SU(3)$ 的结构常数。

	\hspace*{\fill}
	\begin{figure}[htbp]
		\centering
		\includegraphics[width=2.9in]{alpha_s}
		\caption{耦合常数 $\alpha_s$ 随动量转移的演化规律\parencite{Bethke:2006ac}}
		\label{fig:alpha_s}
	\end{figure}

	QCD 理论的基础是夸克和胶子，它们被紧紧地束缚在强子内部，不能够以自由的形态存在于自然界中。QCD 理论中除了夸克和胶子具有与 QED 中电子和光子类似的相互作用以外，胶子之问还存在三胶子和四胶子相互作用顶点。这些顶点使得相互作用耦合强度随着能量的增加而减小~\parencite{Marciano:1977su}。强相互作用的有效耦合常数 $\alpha_s$ 满足：
	\begin{align}
		\alpha_s\left(Q^2\right) \equiv \frac{g^2}{4 \pi}=\frac{4 \pi}{\beta_1 \ln \frac{Q^2}{\Lambda^2}}
		\label{eq:alpha_s}
	\end{align}
	当$g$足够小时，可以用微扰理论来计算$\beta$函数，$\beta_0=(11N_c - 2N_f)/12\pi$，$\beta_1=(11N_c - 2N_f)/3$。$N_c$和$f_f$分别为夸克的色自由度和味道数，标准模型中 $N_c=3$ 而$N_f$取值依赖于标度 $Q$，但$N_f \le 8$即可保证$\beta_0$和$\beta_1$ 都为正直。QCD能量标度为$\Lambda_{QCD} \sim 200 - 300$ MeV，具体值依赖于减除方案和活动味道书。
	
	图~\ref{fig:alpha_s}即是耦合常数 $\alpha_s$ 与动量转移$Q$的关系图。可以看到，耦合常数数 $\alpha_s$ 随$Q$的增大而减小。
	当$Q$的值 趋于无穷大时，即在高能量标度或短距离下，强相互作用的耦合常数$\alpha_s$会趋近于零，此时夸克和胶子的相互作用变得非常弱，即QCD的“渐近自由（asymptotic freedom）”性质，其相互作用可以通过微扰计算来描述~\parencite{Collins:1974ky,Altarelli:1977zs,Politzer:1974fr,Busza:2018rrf}。
	当$Q \to \Lambda_{QCD}$ 时，即在低能下，$\alpha_s$ 变得很大，此时微扰理论不再适用，夸克和胶子束缚在极大的强相互作用力内，此时对应QCD的“色禁闭（confinment）”性质。
	也因此，我们称 $\alpha_s$ 为跑动耦合常数。
	
		%夸克之间越接近，强作用力越弱。当夸克之间非常接近时，强作用力是如此之弱，以便到它们完全可以作为自由粒子活动。这种现象称为“渐近自由”。反之，夸克之间距离越大，强作用力就越强。“渐近自由”的发现导致了全新的理论——量子色动力学（QCD）的诞生。

  \iffalse
  量子色动力学是描述强相互作用的规范理论。QCD的拉氏量函数为:
  \begin{align}
    L=\bar{\psi}(i \not D-m) \psi-\frac{1}{2} \operatorname{tr}\left(F^{\mu v} F_{\mu \nu}\right)
  \label{eq:lagrane}
  \end{align}
  协变微商$D_{\mu}$和
  \begin{align}
    & D_\mu=\partial_\mu-i g T^a A_\mu^a, \\
    & F_{\mu \nu}^a=\partial_\mu A_\nu^a-\partial_\nu A_\mu^a+g f^{a b c} A_\mu^b A_\nu^c,
    \label{eq:Fmunu}
  \end{align}
  $A_{nu}^{a}$ 是相应于规范群 $SU(3)$ 引入的规范场。$g$ 标记着费米子场和规范场 $A_{nu}^{a}$ 之间的相互作用强度。
  \fi
	

 	量子色动力学的发展得到了实验观测的支持。例如，研究者们利用高能粒子对撞实验，观测到了QGP的形成。
 	此外，量子色动力学也与高能核物理、宇宙学等领域相结合，为我们深入理解宇宙的演化和极端条件下的物理过程提供了重要的理论基础。

%------------------------------- 1.2

\section{相对论重离子碰撞和夸克胶子等离子体}

	根据现代的宇宙演化理论，宇宙开始于一次奇异的爆炸（singular explosion），大爆炸爆发出巨大能量，在之后的极端高温环境下，当能量密度达到足够高的水平时，夸克和胶子之间的相互作用变得更加强烈，原本束缚夸克的胶子场开始解禁闭，从而形成一种新的物质状态，即QGP。这种物质的温度如此高以至于质子、中子和其他强子预计会被 "融化"。此时，整个宇宙就可能就处于这种状态中~\parencite{dEnterria:2006mtd}。
	在此之后是一阵膨胀性的扩张，宇宙开始冷却，经历了一系列的相变之后，形成了恒星、星系和地球上的生命，持续到了我们这个时代。

	\hspace*{\fill}
\begin{figure}[htbp]
	\centering
	\includegraphics[width=2.9in]{Phases}
	\caption{有限重子化学势下 QCD 相结构示意图。实线表示相的边界，实心圆表示临界点。横坐标是重子数化学势，纵坐标是温度~\parencite{796947}。}
	\label{fig:Phases}
\end{figure}

	在1974年，李政道就提出可以通过强子对撞机进行高能重离子实验来认为产生QGP所需条件。这一想法引起了领域内物理学家广泛的兴趣。
	如图~\ref{fig:Phases}所示，当温度或密度升高达到相变的临界条件时，夸克和胶子就可以接触禁闭。确定相变临界温度$T_c$以及不同的热力学量在格点量子色动力学（lattice QCD）有所研究。
	在实验上，通过将两个重原子核加速到到接近光速以非常高的能量进行相对论性的碰撞，从而创造了一个极端的高温度、高密度和高压的环境。碰撞区域有可有达到相变的临界条件，夸克和胶子从质子和中子中释放出来，从而实现 QGP 的形成和解禁闭过程~\parencite{Politzer:1973fx,Freedman:1976xs,Shuryak:1977ut}。
	实验上以此来重现类似于宇宙大爆炸后几微秒存在的条件。
	 相对论重离子对撞机的建造就是为了在实验室环境中前所未有的能量密度下创造和研究强相互作用的物质。
	一些研究重离子碰撞的实验设施包括美国布鲁克海文国家实验室的相对论重离子对撞机（Relativistic Heavy Ion Collider，简称 RHIC）和欧洲核子研究中心的大型强子对撞机（Large Hadron Collider，简称 LHC）。
	LHC 在2008年9月10日首次启动，其重离子项目的主要物理学目标是产生和测量QGP的特性。对撞机内的光束在加速器环周围的四个位置进行碰撞，对应于四个粒子探测器的位置 —— ATLAS、CMS、ALICE和LHCb。
	随着实验设备的不断完善，目前，LHC实验上对于铅-铅碰撞上升到了更高能及，并且氙-氙碰撞也有了新的结果~\parencite{CMS:2023pad,Feuillard:2023xgh,Stanek-Maslouska:2023xrq,ATLAS:2023vms,ATLAS:2023xzy}。
	高能核碰撞作为探索基本粒子性质和理解宇宙演化的重要手段，在过去几十年里取得了显著的进展。
	这些实验这提供了研究碰撞中产生的物质的独特方法，导致了许多新粒子的发现，如希格斯玻色子。
	对其中硬散射过程的研究可以帮助研究人员探测QGP的特性，以及控制夸克和胶子之间相互作用的强作用力。
	在这些碰撞中产生的高能粒子可以被探测和分析，以更好地了解在这些极端条件下发生的基本相互作用和动力学。
	下一节将介绍 QGP 的形成和演化过程。

%------------------------------- 1.2.1
\subsection{重离子碰撞的时空演化}

	在相对论重离子碰撞中经历了相变过程，夸克和胶子逐渐重新组合成重子和介子，而强相互作用、色荷传播和强子化学势的变化影响了粒子的产生、流动和相互作用。
	通过对其碰撞的时空图像的研究，可以深入理解QGP的性质、相变过程以及强子产生和相互作用的动力学。这对于深入了解宇宙早期的宇宙学条件、量子色动力学的非扰动效应以及强子物质的性质具有重要意义。
	图~\ref{fig:HIC}展示了两个接近光速运动的原子核在时间 $t=0$ 和纵向位置 $z=0$（未显示横向方向）发生碰撞，随后经历了不同的阶段。火球除了纵向膨胀外，还在径向方向上膨胀。相对论重离子碰撞的时间演化可以大致分解为三个明显的阶段：初始阶段、夸克-胶子等离子体（QGP）阶段和强子气相阶段。下面作详细描述：

	\hspace*{\fill}
	\begin{figure}[htbp]
		\centering
		\includegraphics[width=3.1in]{HIC}
		\caption{重离子碰撞的光锥演化示意图。演化过程从前平衡阶段的热火球开始（灰色区域）、通过形成夸克-胶子等离子体阶段（红色区域）随后经历了一个交叉（cross-over）相变（黄色区域）到强子气体的阶段（蓝色）。化学冻结和动力学冻结在强子相也有所显示~\parencite{Braun-Munzinger:2018hat}。}
		\label{fig:HIC}
	\end{figure}


	\begin{itemize}
		\item 碰撞的初始状态：
		在实验中，两束高能重离子（如铅离子）以接近光速相撞，随后相互穿透时，参与粒子经历了核子-核子的非弹性散射和初态辐射，形成大量高度激发的部分子。大量能量被沉积到碰撞区域中。
		新形成的介质处于预平衡（pre-equilibrium）阶段，并通过介质成分之间的频繁碰撞接近局域热平衡。在这个阶段，还通过大$x$散射产生了被称为喷注（jet）的高能量粒子。随机的碰撞参数$b$在重叠区域的横向平面上产生涨落的离心率。理解初始几何形状对解释碰撞的实验结果很有必要。
		
		\item 热化过程以及夸克胶子等离子的体的形成：
		系统的温度足够高时，夸克和胶子的相互作用变得非常强烈，高能量密度使得核子分解成了退禁闭的夸克和胶子。夸克和胶子由束缚状态变成自由状态，逐渐充满了整个碰撞区域，
		%当火球经历适当时间 $\tau =  (\tau^2 − z^2)^{1/2}  \sim 1$  fm/$c$后，会接近于局部热平衡状态，
		形成一个高温高能量的夸克胶子等离子体（QGP）。
		%在该状态下，QGP被认为是一种几乎理想的流体，具有非常低的切向粘滞与熵密度比（$\eta/s$），其演化在相对论流体力学的框架下得到很好的描述。
		这个解紧闭物质存在的时间很短，在LHC上大约为$10-15$ fm$/c$。
		%火球在LHC上的温度超过了 $300$ MeV~\parencite{ALICE:2015xmh}。
		QGP的性质和演化过程受到量子色动力学的影响。%在这个等离子体中，夸克和胶子之间的相互作用仍然存在，但由于高温和高能量密度的特殊条件，它们表现出与常规核物质不同的行为。
		由于介质的高密度和平均自由程相对于系统尺度而言较小，介质中的部分子很快建立了热平衡并表现出集体行为。
		QGP的集体运动可以很好地用相对论流体动力学描述，其中使用了QCD 状态方程作为计算的输入，即在给定温度下压强、能量密度和强子密度之间的关系~\parencite{Heinz:2013wva,Heinz:2004qz,Jeon:2015dfa,Kolb:2003dz,Nair:2006qg,Huovinen:2003fa}。
		%关于相对论流体力学在（极端）相对论重离子碰撞中的应用可以参考这几篇综述文章。
		% 由于QGP是色不透明的，碰撞早期形成的彩色硬探针（如喷注）在穿过 QGP 时会经历能量损失，即喷注淬火。对余下喷注的测量提供了QGP的性质。
		随着系统的演化，热压向周围真空施加的力使得 QGP 经历集体膨胀和冷却。
		
		\item 强子相的冻结：
		随着时间的推移，QGP开始迅速膨胀并冷却，当低于临界温度$T_c$时，开始发生强子化反应，热化的部分子开始重新组合形成色中性的强子，如质子和中子。
		% 即使在QGP转变为强子物质之后，粒子仍然可能相互作用和散射，保持系统的局部平衡，尽管在强子化阶段的膨胀过程中，密度降低到足够低，使得非弹性碰撞很少发生。与此同时，处于激发态的不稳定强子会衰变为稳定的强子，这被称为共振衰变。
		% 当系统的温度（局部地）达到Tc≈化学冻结温度Tchem时，
		一旦所有强子都处于稳定状态，粒子的产生数目就会定格。%尽管弹性碰撞可能仍会发生，对应的温度被称为化学冻结温度$T_{chem}$。%在动力学冻结之前，强子之间的再散射仍然存在。
		% 在动力学冻结之后，物质种类和能量-动量保持不变，直到在探测器中被探测到。
		随后可以通过探测器来观测和记录，以分析出有关夸克胶子等离子体性质的信息。
		%强子的产量可以被确定并被"冻结"。随后在强子相中的演化只能通过非平衡方法来描述~\parencite{Xu:2017akx}。 通常情况下，它被建模为一个强子级联模型，其中进一步的弹性和原则上的非弹性散射可能发生。热力学模型对于T = Tchem的预测与实验数据的出色一致性暗示非弹性再散射应该很小。。粒子在动力学冻结温度Tkin以下停止散射并自由流动。相对论核碰撞中粒子产生研究的一个重要结果是，对于在RHIC和LHC上研究的质心能量，Tchem的值与Tc非常接近：化学冻结发生在QCD相边界附近。	
		%事实上，z方向（束流方向）的膨胀是这些碰撞中的主要冷却机制，这个一维时空图可以完全通过解析方法求解（Bjorken模型~\parencite{Bjorken:1982qr}）。火球演化的不同阶段通常通过不同的模型方法来处理。例如，前平衡阶段（$\tau_0 \le 1$ fm/c）通常使用部分子级联（parton-cascade）模型或玻尔兹曼输运（Boltzmann transport）方法~\parencite{Heinz:2004qz,Heinz:2015arc,Liu:2015nwa,Braun-Munzinger:2015hba}进行量子输运建模。
	\end{itemize}

	\hspace*{\fill}

%------------------------------- 1.2.2

\subsection{夸克胶子等离子体形成的信号}
	
	  相对论重离子物理学的一个主要研究对象是QGP，而它的存在时间极端，约为 10fm，因此从实验上无法直接对其测量。我们想要研究QGP就得借助一些能够探测到它的信号，即QGP探针。以下将介绍一些QGP探针。

	   1. 直接光子和双轻子的产生：
	重离子碰撞中产生的大量夸克和胶子发生相互作用，其中夸克和反夸克发生散射产生胶子和光子$q+\bar{q} \to g+\gamma$，夸克（或反夸克）与胶子也会发生散射产生光子$q(\bar{q}+g\to q(\bar{q}+\gamma)$，这些热光子从QGP中产生但并不参与强相互作用，因此携带着QGP的一些信息，可以用来提取出QGP中各种热力学参数 ~\parencite{McLerran:1984ay,Rapp:1999ej}。
	  
	  2. 奇异粒子的增加：
	  奇异夸克（$s$）不是组成物质的原始粒子，但是QGP中存在有大量的夸克和胶子，他们之间相互作用会产生大量的$s\bar{s}$，之后会强子化形成大量的奇异强子，导致末态奇异强子产额增加，这被认为QGP存在的信号。奇异强子的产生提供了QGP温度，化学势等信息~\parencite{Rafelski:1982pu,Soff:1999et,Koch:1986ud}。
	  
	  3. $J/\psi$ 产额的压低：
	   $J/\psi$ 是粲夸克（$c$）与其反夸克（$\bar{c}$）的束缚态，即粲夸克偶素（$c\bar{c}$）。它产生于QGP形成之前，当QGP温度足够高时，德拜屏蔽效应减弱会使得$J/\psi$解离而再度成为自由的$c$和$\bar{c}$，于是在实验上就观察到了$J/\psi$ 产额的压低，这也是目前重味物理研究的主要课题之一
	  ~\parencite{Brambilla:2010cs,Richardson:1978bt,Thews:2000rj}。
	  
	\hspace*{\fill}
\begin{figure}[htbp]
	\centering
	\hspace{-0.5cm}
	\includegraphics[width=4in]{anisotropy}
	\caption{初始重叠区域示意图（左）；由这个各向异性重叠区域产生的空间各向异性的示意图（右）。其中时间平面为垂直于入射粒子束（$z$轴）方向的平面~\parencite{Connors:2017ptx}。}
	\label{fig:anisotropy}
\end{figure}
	  
	  4. 集体流效应：
	  在高能重离子碰撞中，产生了具有相当高的密度和强相互作用的介质，因此系统很快就会实现热平衡，介质产生的热力学压力导致系统的集体膨胀。初态碰撞的叠加区域会因为入射部分子位置的涨落以及碰撞参数的影响而表现为位置上的不对称性。如图~\ref{fig:anisotropy}所示，这种不对称性表现为压强梯度沿$y$轴方向较小且沿$x$轴方向较大。经历预平衡后热密介质会因此沿$y$轴方向膨胀得更快，即“各向异性”，这会随强子化最终体现在末态强子的动量分布上。我们可以使用动量各向异性的傅里叶系数来量化这种各向异性：
\begin{equation}
	E \frac{d^3 N}{d^3 p}=\frac{1}{2 \pi} \frac{d^2 N}{p_T d p_T d y}\left(1+2 \sum_{n=1}^{\infty} v_n\left(p_T, y\right) \cos \left[n\left(\phi-\psi_R\right)\right]\right)
\end{equation}
	其中，$v_n$为傅里叶展开的$n$阶系数，$n= 1, 2, 3$分别表示直接流、椭圆流和三角流，$\phi$和$\psi_R$分别为末态强子和反应平面的方位角。
	目前实验上通过测量最多的是末态强子的椭圆流信号，观察到的集体行为成为QGP形成的标志。
	~\parencite{Ollitrault:1992bk,Heinz:2013th,PHENIX:2003qra,PHENIX:2006iih,ALICE:2013olq}。
	  

	  
	  5. 喷注淬火效应：
	  在极端相对论重离子碰撞中，由高能夸克和胶子引起的喷注（jet）也能为QGP的存在提供一种探针。
	  喷注是携带着与散射部分子的能量大致相当的一束校准的粒子束。这些粒子或它们的能量被聚类在一起形成被测量的喷注。
	  来自初始的高能硬碰撞的部分子通过在碰撞后形成的核之间的密集介质中传播并相互作用，导致能量损失，从而影响它们的产额谱。
	  预计在核-核碰撞中的能量损失相对于质子-质子碰撞或质子-核碰撞中更大，这一现象被称为喷注淬火（jet	 quenching）~\parencite{Gyulassy:2003mc,Wang:1992qdg,Mehtar-Tani:2013pia,Blaizot:2015lma,Qin:2015srf,JET:2013cls,Majumder:2010qh,Wang:1991hta,Gyulassy:1990ye,Zapp:2012ak,ALICE:2019qyj,ATLAS:2018gwx}。也就是说，硬喷注与软介质之间的相互作用通常会减少传播中的喷注部分子的能量。这会导致喷注内部的高能粒子减少，能量分布变得更加宽泛，甚至出现分裂和消失的现象~\parencite{Qin:2015srf,Wang:2001cs}。 
	  这最早在重离子碰撞中的RHIC和LHC上进行了一些测量~\parencite{CMS:2016uxf,CMS:2011iwn,STAR:2002svs,PHENIX:2001hpc}。 

	量化介质效应对高能粒子产额谱的影响可以用“核修正因子（nuclear modification factor）”来描述~\parencite{ STAR:2003pjh}：
	   \begin{equation}
	  	R_{AB}(p_T)=\frac{d \sigma^{A B}_h /d p_T}{N_{col}^{A B} ~ d \sigma^{p p}_h / d p_T}
	  \end{equation}
  
    \begin{equation}
  	R_{\mathrm{AA}}=\frac{1}{N_{\mathrm{evt}}} \frac{d N_{\mathrm{jet}}}{d p_{\mathrm{T}}} /\left(\left\langle T_{\mathrm{AA}}\right\rangle \frac{d \sigma_{p p}}{d p_{\mathrm{T}}}\right)
  	\label{eq:RAA}
  \end{equation}

	  其中$d \sigma^{A B}_h /d p_T$和$d \sigma^{p p}_h / d p_T$分别为核-核碰撞和质子-质子碰撞中的横动量分布，$N_{col}^{A B}$为核子-核子两体碰撞次数。
	  即核-核碰撞中的产额除以参与碰撞的核子对数 $N_{b i n}^{A B}(b)=\int d^2 r t_A(r) t_B(|\vec{b}-\vec{r}|)$，此结果再与相同条件下核子-核子碰撞的产额比较。
	  核修正因子用于将实验测量得到的高能粒子数谱在核-核碰撞中与预期的参考谱进行比较。
		核修正因子对1的偏离体现了碰撞中形成的热密核物质的介质效应。
	  当核修正因子小于1时，说明介质效应导致高能粒子产额谱的减少，即介质压低效应；在没有核效应的情况下，核-核碰撞是核子-核子碰撞的简单叠加，核修正因子为1。
	  RHIC的STAR实验测量了在质心系能量$\sqrt{200}GeV$下，首次观察到了在金-金碰撞的中心区域，高横动量下强子产额的压低见图~\ref{eq:RAA1}，即所谓的喷注淬火现象，原因即是高横动量的部分子在穿过QGP时与其发生相互作用而损失能量。
\begin{figure}[htbp]
	\centering
	\hspace{-0.5cm}
	\includegraphics[width=3.5in]{RAA1}
	\caption{在质心系能量（centre-of-mass energy）下$\sqrt{200}GeV$ 金-金对心碰撞和氘-金碰撞的核修正因子分布~\parencite{STAR:2003pjh}。}
	\label{fig:RAA1}
\end{figure}
	近年来，随着LHC和RHIC提供的越来越精确的测量数据以及理论的不断改进，我们对于重离子碰撞中喷注淬火的理解发生了巨大的变化。此外，用于测量喷注子结构的技术不断提升。目前的关注点是理解喷注淬火与部子簇射结构以及喷注穿越QGP的长度之间的关系。除了对喷注本身的修正和淬火进行研究外，研究QGP对喷注穿越过程的响应也具有重大意义。



在重离子碰撞中，对整体喷注（full jet）及其在核物质中的核修正的研究与单强子产生的研究有着显著区别，同时也预期能够提供更详细的喷注与介质相互作用信息。在使用单个非遍举强子进行喷注淬火的研究时，主要关注点在于领头强子与介质组分的相互作用。
在这种情况下，与介质组分的辐射非弹性相互作用通常被认为是负责部分子能量损失和重离子碰撞中轻夸克强子产生压低的主要机制~\parencite{Qin:2007rn}，而弹性碰撞通常为轻味喷注淬火提供了补充贡献。需要注意的是，对于重味，由于有质量效应的存在，特别是在低横动量下~\parencite{Cao:2013ita}，碰撞过程变得更加重要。







%------------------------------- 1.3

\section{双喷注物理的研究现状}
%\subsection{高能重离子碰撞中的开放重味物理研究现状***}

	在高能核碰撞中，双喷注物理的研究是粒子物理学的重要领域之一。双喷注（Dijet）是指初始部分子的QCD演化产生的背对背喷注，即成对的高能喷注。这些喷注通常由强子碎裂产生，其横向动量较大。
	研究人员利用高能核碰撞实验中观测到的双喷注事件，可以探索量子色动力学在高能量下的特性并验证和预测微扰QCD的性质。

	在20世纪60年代，深度非弹性散射（deep inelastic scattering，简称DIS）实验显示出质子和中子具有内部结构：夸克和胶子。从基本组分粒子的角度来理解质子的构成是所有电子-质子散射实验的主要目标之一。
	在绘制核子组分的一维动量分布方面，以费曼部分子分布函数（parton distribution functions，简称PDF）的形式，取得了巨大的进展。
	然而，后续的实验数据表明自由核子的图像过于简单~\parencite{EuropeanMuon:1981gyc,Bodek:1983qn}。核子中的夸克和胶子对其核环境非常敏感，因此当将自由核子在重核中时，质子的部分子分布会发生变化，因为来自相邻核子的夸克和胶子非常接近，例如来自不同核子的两个胶子可以融合在一起。
	也就是说，原子核中的部分子分布函数（Nuclear parton distribution functions，简称nPDF）与自由核子的是不同的。在固定靶的DIS实验数据和$p$A碰撞的Drell-Yan实验的产生中，nPDFs可以通过对广泛的实验数据进行全局量子色动力学分析来被确定。理解夸克和胶子在核环境中的行为是现代核物理的一个重要研究重点。	这些研究不仅揭示了核子的部分子结构，而且为研究强相互作用提供了重要机会。
%---------------
	在过去几年中，LHC上的质子-铅碰撞提供了丰富的机会来研究冷核物质效应~\parencite{ATLAS:2015mwq,ALICE:2016rzo,ALICE:2020jff,CMS:2021ynu,LHCb:2022kph,CMS:2019leu,ALICE:2022cxs}。
	特别是，物理学家通过各种努力从高能实验中确定夸克和胶子的PDFs~\parencite{Hou:2019efy,Cocuzza:2021rfn,NNPDF:2017mvq,Alekhin:2017kpj}。
	喷注大多是在$2 \to 2$ 的散射过程中成对产生的。
	具有较大横向动量的双喷注事件可以在QCD的背景下描述为部分子-部分子散射。
	双喷注截面的测量对于验证微扰QCD的准确性至关重要。通过比较实验测量结果与理论预测，可以检验QCD的基本假设和计算方法。此外，双喷注截面的测量还可以提供有关部分子分布函数（PDFs）的信息。
	Dijet 方位角的相关性是直接探测硬散射中涉及的胶子动力学的最敏感的方法之一~\parencite{Marquet:2007vb,Dominguez:2011wm}。
	在LHC的质子-铅碰撞中，也研究了$p$A碰撞中的双喷注的相关性，对小$x$物理做了运动学范围的补充~\parencite{ATLAS:2019jgo}。而在理论的解释中，结合了胶子饱和和高阶软胶子辐射的贡献~\parencite{vanHameren:2019ysa}。
	此外，双喷注事件也在实验上通过横能$E_T$、快度$y_1, ~y_2$研究了三重微分~\parencite{CDF:1989car}，也通过平均横向动量及$p_{T,avg}$、快度变量$dy_b$和$dy^*$研究了三重微分~\parencite{CMS:2017jfq}。
	这些实验结果对于验证QCD理论、改进模型和理解强子的强相互作用起着重要作用。


%---------------
%---------------
	在重离子碰撞的初始阶段，入射核的组分（质子和中子）之间会发生硬散射相互作用。
	这些碰撞涉及大动量转移，导致高能粒子的产生。
	当喷注部分子穿过夸克胶子等离子体介质时，与介质之间的强相互作用，经历了弹性散射~\parencite{bjorken1982fermilab,braaten1991energy,Djordjevic:2006tw,Qin:2007rn}和非弹性散射~\parencite{Baier:1996kr,Zakharov:1996fv,Gyulassy:1999zd,Wiedemann:2000za,Arnold:2002ja,Wang:2001ifa}。
	这些相互作用不仅导致了能量-动量在介质中的沉积和从介质中的抽取（“喷注诱导的介质激发”或“介质响应”），还会导致喷注的核修正。
	在过去的二十年中，为了揭示这些部分子的强相互作用效应，研究人员发展了一系列的方法。这些工具包括高横动量强子或喷注的压低因子$R_{AA}$ ~\parencite{ATLAS:2014ipv,Wang:1998bha,Bass:2008rv,Wang:2002ri,JETSCAPE:2021ehl,He:2018xjv}、双喷注的动量不对称性~\parencite{ATLAS:2010isq,Zhang:2007ja,Qin:2010mn,Kang:2018wrs,Fialkovsky:2011pu}、矢量玻色子与喷注的相关性（$\gamma/Z^0 +$jet） ~\parencite{CMS:2012ytf,Neufeld:2010fj,Dai:2012am,Huang:2015mva,Yang:2022nei,Chen:2020tbl}、全局事件几何形状 ~\parencite{Chen:2020pfa,Kang:2023qxb}和喷注子结构 ~\parencite{CMS:2013lhm,Vitev:2009rd,Vitev:2008rz,Zhang:2021sua,Chen:2022kic,Zhang:2018urd,Wang:2023eer,JETSCAPE:2023hqn}。
	研究这些性质，为我们理解这种强耦合夸克物质的性质以及在极高温高密度条件下测验量子色动力学这一基本理论开辟了新的途径~\parencite{Xie:2022ght,Zhang:2021xib,JETSCAPE:2020mzn,Tang:2020ame}。
%---------------
	与单喷注的压低相比，Pb+Pb对撞中的双喷注为研究喷注淬火提供了一个更干净的探针。由于在QGP中穿越了不等长的路径，这两个喷注就会经历不对称的能量损失，这是由重叠核的几何形状以及穿过不断变化的介质的喷注轨迹的相对方向所驱动的~\parencite{Qin:2010mn}。
	 其中在对$R_{AA}$的测量中可以看到，与$pp$对撞中的产额相比，Pb+Pb中心碰撞中的非遍举喷注的产额被压低了2倍。虽然对非遍举喷注产额压低的测量包含了喷注在QGP中穿越不同路径的长度和受到不同的喷注淬火波动(fluctuations)的影响的贡献，但对双喷注的测量提供了关于能量损失的路径长度依赖性和能量损失中的波动作用的额外信息~\parencite{Qin:2010mn, Milhano:2015mng}。
	铅-铅碰撞中的双喷注压低是用动量不平衡(momentum imbalance)~\parencite{CMS:2012ulu, ATLAS:2017xfa, ATLAS:2022zbu}和对核修正因子~\parencite{ATLAS:2022zbu}来衡量的，它量化了领头喷注（在横动量方向）和次领头喷注（在方位角相反方向）压低之间的差异。
	到目前为止，只在XeXe碰撞中测量了非遍举带电强子的压低效应，而涉及喷注和双喷注的部分子能量损失测量仅在PbPb碰撞中进行。
	2017年LHC XeXe运行为研究比Pb更轻的核碰撞中的喷注淬火提供了可能性，这是非常有吸引力的。在重离子碰撞的最中心碰撞中，即碰撞几何形状是最对称时，XeXe碰撞中的基本事件（Underlying Event）比在Pb+Pb碰撞中的小~\parencite{ATLAS:2022kqu, CMS:2018yyx, ALICE:2018hza}；其次，铅核和氙核之间的核子数量或核半径的减少可能会影响喷注淬火的程度，因为在介质中硬散射的部分子所穿过的路径长度和总能量密度都会减少。
	此外，在不同的碰撞系统中对双喷注的测量会进一步约束路径长度、能量密度和对喷注淬火的波动的影响。
	而在LHC中涉及喷注和双喷注的部分子能量损失的测量只在Pb+Pb对撞中有进行。



%------------------------------- 1.6
\section{本文提纲}

		在前文中，笔者介绍了标准模型的理论框架以及描述强相互作用的理论，即量子色动力学。同时也介绍了相对论重离子碰撞形成了一个非常短暂的高能密度区域，产生了一种新的物态，即夸克胶子等离子体。
	% Page 12
		我们工作的出发点是利用$p$Pb碰撞中的双喷注产生，来研究部分子分布函数及其核修正，特别是胶子的分布。我们研究了4种不同的三重微分散射截面，目的是在观测量层面上，较好的区分动量分数$x$和标度$Q$。
	我们希望在领头阶，
	%可以做到像DIS实验中的那样，
		尽可能的减小高阶修正。
		随后，我们研究了双喷注的压低和关联，第一次从理论角度研究了XeXe碰撞中质心能量为$\sqrt{s}=5.44$~TeV下的双喷注动量不平衡量$x_J$，相关计算并同ATLAS实验符合较好。此外对$x_J$的介质修正也进行了研究，并探讨了领头喷注和次领头喷注之间横向动量的不平衡性。
		对双喷注的初始产生，我们采用了POWHEG+PYTHIA8的方案，该方案将次领头阶的QCD矩阵元与部分子簇射效应相匹配。双喷注的介质演化由输运模型描述，该方法考虑了QGP中的弹性和非弹性部分子的相互作用。
		此外，我们以非遍举、$c\bar{c}$和$b\bar{b}$三种双喷注的$\Delta x_J$分布研究了质量效应。
		我们还计算了$\sqrt{s}=5.44$~TeV的XeXe碰撞和$\sqrt{s}=5.02$~TeV的PbPb碰撞的双喷注的核修正因子比值，且与近期的ATLAS实验结果也进行了比较
	本文余下内容分为四部分：
	在第二章中，
		从理论物理的角度着重介绍了高能核物理的基本理论与研究方法，比如利用微扰展开来计算物理过程的各阶贡献、几种用于识别和重构喷注的算法、以生成模拟数据便于与实验数据进行比较的蒙特卡洛事件产生器、计算碰撞事件中的各种几何特征的Glauber模型、核子内部的部分子携带动量分布的函数及其核修正、提供介质中粒子的演化和传播信息的输运模型。
	在第三章中，
		我们主要研究了束缚核子相对于自由核子的修正效应（表示为$r$），并发现不同核效应之间存在着显著差异，尤其是对于胶子分布。。
		我们研究了作为QGP硬探针基准的冷核物质效应，以双重微分散射截面的计算为例，
		讨论了$p$A喷注中双喷注的产生机制，使用了四种核修正的参数化形式：EPPS16、nCTEQ15、TUJU19 和 nIMParton16，发现他们之间在观测量层面上也存在差异。
		我们探讨了通过末态观测量来提取初态信息的可行性，以此解释不同类型的部分子分布函数。
	在第四章中，
		我们提出研究双喷注的三重微分散射截面的重要性，讨论了在$p$A 碰撞中这种截面对于揭示PDF核修正方面有更多优势，而且相应观测量可以作为 $r^A_i(x,Q^2)$ 从小到大 $x$ 值的良好成像。
		可以通过同时固定三个双喷注的变量，建立与 $x$ 和 $Q^2$ 的更好的对应关系。
	在第五章中，
	


	最后在第六章中做了总结和展望。
	最后阐明本论文采用了自然单位制：$\hbar = c = k_B = 1$，其中$c$为光速，$\hbar = h/2\pi$，$h$为普朗克常数。
	常规张量$g^{\mu\nu} = diag(1,-1,-1,-1)$。


